\input{header.tex}
\input{newcommands.tex}
% ===================================================================
%  Vorlage: Mattaufgabe „Matt in N“ mit Schritt-für-Schritt-Lösung
%  Seite 1: Anfangsstellung und Aufgabe, danach je Zug ein Brett mit
%  Pfeilen und Erklärung. Beispiel: Matt in 11 (ersticktes Matt).
%  Brett farbig über \schachbrettfarben (siehe references/befehle.md).
%
%  Für eine eigene Aufgabe anpassen: \StartFEN, die Zugfolge in
%  \hidemoves und \mainline, die \Schritt-Zeilen und die Texte.
%  Stellung vorher prüfen:  python3 pruefe_stellung.py --matt N "FEN"
%  (Bei langen Matts wie hier zusätzlich mit einer Engine prüfen.)
%
%  Quelle der Beispielstellung: YouTube-Short „Insane Knight Checkmate“,
%  https://www.youtube.com/shorts/F65ssb7mKkE
%  Kompilieren: pdflatex vorlage_matt_in_n_loesung.tex
%  (header.tex und newcommands.tex aus latex-grundstil daneben legen)
% ===================================================================

% --- Brettfarben: Standard ist blau (aus der newcommands.tex). Andere Schemata:
%     \schachbrettfarben{braun}, {gruen}, {grau} oder {klassisch} (schraffiert).
%\schachbrettfarben{braun}

\pagestyle{plain}
\setstretch{1.15}
\setlength{\parskip}{0.6em}

% --- Die Stellung (einmal zentral, alles andere greift darauf zu)
\newcommand\StartFEN{3N1qk1/6pp/5p2/8/1R2R3/K1p5/6Q1/4r2r w - - 0 1}

% \Schritt{Zug-ID}{Pfeilstil Weiß}{von}{nach}{Antwort Schwarz}{Überschrift}{Erklärung}
% Das Brett zeigt die Stellung NACH dem weißen Zug (aus der gespeicherten Partie geholt).
% Gelbe Felder + grüner Pfeil = weißer Zug (straightmove oder knightmove),
% roter Pfeil = Antwort von Schwarz, die danach folgt.
\newcommand\Schritt[7]{%
	\par\noindent
	\begin{minipage}[c]{0.33\linewidth}\centering
		\xskakset{id=matt, moveid=#1}%
		\chessboard[setfen=\xskakget{nextfen}, smallboard, showmover=false,
			pgfstyle=color, color=brettmarke, colorbackfields={#3,#4},
			linewidth=0.09em, shortenend=0.35em,
			pgfstyle=#2, color=akzentgruen, markmoves={#3-#4},
			pgfstyle=straightmove, color=akzentrot, markmoves={#5}]
	\end{minipage}\hfill
	\begin{minipage}[c]{0.63\linewidth}
		{\large\bfseries #6}\par\smallskip
		#7
	\end{minipage}\par\bigskip}

% Figurensymbole im Fließtext
\newcommand\K{\symking}
\newcommand\D{\symqueen}
\newcommand\T{\symrook}
\newcommand\Sp{\symknight}


\begin{document}

% Partie einmal unsichtbar durchspielen – skak prüft dabei jeden Zug
\newchessgame[id=matt, setfen=\StartFEN]
\hidemoves{1. Qa2+ Kh8 2. Nf7+ Kg8 3. Nd6+ Kh8 4. Rb8 Qxb8 5. Nf7+ Kg8
	6. Nd8+ Kh8 7. Re8+ Rxe8 8. Nf7+ Kg8 9. Nh6+ Kh8 10. Qg8+ Rxg8 11. Nf7#}

%%%%%%%%%%%%%%%%%%%%%%%%%%%%%%%%%%%%%%%%%%%%%%%%%%%%%%%%%%%%%%%%%%%%%%
%  SEITE 1: AUFGABE
%%%%%%%%%%%%%%%%%%%%%%%%%%%%%%%%%%%%%%%%%%%%%%%%%%%%%%%%%%%%%%%%%%%%%%
\begin{center}
	{\LARGE\bfseries\sffamily Das Springer-Karussell}\par\smallskip
	{\large Eine Mattaufgabe mit zwei Turmopfern und einem Damenopfer}
\end{center}

\bigskip
\begin{center}
	\chessboard[setfen=\StartFEN, normalboard, showmover=true]
\end{center}

\begin{center}
	\begin{tabular}{@{}ll@{}}
		\textbf{Weiß:}    & \K a3, \D g2, \T b4, \T e4, \Sp d8 \\
		\textbf{Schwarz:} & \K g8, \D f8, \T e1, \T h1, \sympawn c3, \sympawn f6, \sympawn g7, \sympawn h7 \\
	\end{tabular}
\end{center}

\medskip
\schwer \textbf{Aufgabe.} Weiß ist am Zug und setzt in \textbf{11 Zügen} matt.
Finde die vollständige Zugfolge. Fast jeder weiße Zug ist ein Schach, Schwarz hat fast nie eine Wahl.

Weiß muss sich beeilen: Wäre Schwarz am Zug, käme er mit 1\,\ldots\,\T a1+ selbst zu einem
entscheidenden Angriff auf den weißen König. Ruhige Züge kann sich Weiß also (fast) nicht leisten.

\begin{hinweis}[Tipps (erst lesen, wenn du nicht weiterkommst)]
	\begin{enumerate}[1.]
		\item Der erste Zug bringt die Dame mit Schach auf die Diagonale zum schwarzen König.
		\item Am Ende steht ein \textit{ersticktes Matt}: Der König in der Ecke ist von eigenen Figuren
		eingemauert, ein Springer setzt matt.
		\item Solange die schwarze Dame auf f8 steht, klappt das nicht. Sie muss weggelockt werden, und
		am Ende muss ein schwarzer \textit{Turm} auf g8 stehen.
	\end{enumerate}
\end{hinweis}

\vfill
{\footnotesize Stellung als FEN: \texttt{\StartFEN}\par
Quelle der Stellung: YouTube-Short \enquote{Insane Knight Checkmate} (Kanal @sigma\_chess21),
\url{https://www.youtube.com/shorts/F65ssb7mKkE}}

\newpage
%%%%%%%%%%%%%%%%%%%%%%%%%%%%%%%%%%%%%%%%%%%%%%%%%%%%%%%%%%%%%%%%%%%%%%
%  LÖSUNG
%%%%%%%%%%%%%%%%%%%%%%%%%%%%%%%%%%%%%%%%%%%%%%%%%%%%%%%%%%%%%%%%%%%%%%
\section*{Lösung: Matt in 11 Zügen}

\newchessgame[id=druck, setfen=\StartFEN]
\mainline{1. Qa2+ Kh8 2. Nf7+ Kg8 3. Nd6+ Kh8 4. Rb8!! Qxb8 5. Nf7+ Kg8
	6. Nd8+ Kh8 7. Re8+! Rxe8 8. Nf7+ Kg8 9. Nh6+ Kh8 10. Qg8+!! Rxg8 11. Nf7#}

\begin{hinweis}[Die Idee]
	Das bekannte Muster des erstickten Matts lautet: Springerschach auf f7, Doppelschach auf h6,
	Damenopfer auf g8, Springermatt auf f7. Hier scheitert der direkte Weg
	1.\,\D a2+ \K h8 2.\,\Sp f7+ \K g8 3.\,\Sp h6+ \K h8 4.\,\D g8+ aber an 4\,\ldots\,\D xg8:
	Die Dame deckt von g8 aus das Feld f7, und nach 5.\,\Sp f7+ \D xf7 ist der Angriff vorbei.
	Auf g8 muss deshalb ein \textit{Turm} schlagen. Weiß opfert dafür erst einen Turm auf b8
	(lenkt die Dame von f8 weg) und dann den zweiten auf e8 (lockt einen schwarzen Turm auf die 8.~Reihe).
	Der Springer pendelt dabei immer wieder über f7 – sieben der elf weißen Züge macht er allein.
\end{hinweis}

\textbf{So liest du die Bretter:} Jedes Brett zeigt die Stellung \textit{nach} dem weißen Zug.
Die gelben Felder und der \textcolor{akzentgruen}{\textbf{grüne}} Pfeil zeigen den weißen Zug, der
\textcolor{akzentrot}{\textbf{rote}} Pfeil die Antwort von Schwarz, die danach folgt.

\bigskip

\Schritt{1w}{straightmove}{g2}{a2}{g8-h8}
{1.\,\D a2+ \quad \K h8}
{Die Dame gibt auf der Diagonale a2--g8 Schach. Schwarz kann nur in die Ecke ausweichen.
	Dazwischenziehen verliert sofort: 1\,\ldots\,\D f7 2.\,\D xf7+ \K h8 3.\,\D e8\#.}

\Schritt{2w}{knightmove}{d8}{f7}{h8-g8}
{2.\,\Sp f7+ \quad \K g8}
{Der Springer gibt Schach. Schlagen darf Schwarz ihn nicht: Nach 2\,\ldots\,\D xf7 hat die Dame
	die 8.~Reihe verlassen, und es folgt 3.\,\T b8+ \D g8 4.\,\T xg8\#. Also muss der König zurück nach g8.}

\Schritt{3w}{knightmove}{f7}{d6}{g8-h8}
{3.\,\Sp d6+ \quad \K h8}
{\textbf{Abzugsschach:} Der Springer zieht weg und öffnet die Diagonale der Dame.
	Auf d6 steht er genau richtig: Er hat die 8.~Reihe für den Turm freigemacht und behält f7 im Blick.
	Wieder geht nur \K h8, denn nach 3\,\ldots\,\D f7 wäre die Dame gefesselt und 4.\,\T b8 matt.}

\Schritt{4w}{straightmove}{b4}{b8}{f8-b8}
{4.\,\T b8!! \quad \D xb8}
{Das erste Turmopfer und der einzige weiße Zug ohne Schach. Die Dame auf f8 ist angegriffen und zugleich
	\textbf{gefesselt}: Sie darf die 8.~Reihe nicht verlassen, weil dahinter der König steht.
	Es droht \T xf8\#. Schlägt die Dame den Turm, ist sie von f8 weggelenkt und deckt f7 nicht mehr.
	(Andere Züge wie 4\,\ldots\,h6 verlieren nach 5.\,\T xf8+ noch schneller.)}

\Schritt{5w}{knightmove}{d6}{f7}{h8-g8}
{5.\,\Sp f7+ \quad \K g8}
{Der Springer kehrt nach f7 zurück. Jetzt kann ihn niemand mehr schlagen, der König hat nur das Feld g8.}

\Schritt{6w}{knightmove}{f7}{d8}{g8-h8}
{6.\,\Sp d8+ \quad \K h8}
{Wieder Abzugsschach, diesmal zurück nach d8. Dort versperrt der Springer der schwarzen Dame die
	8.~Reihe: Sie kommt von b8 nicht mehr nach e8. Auf 6\,\ldots\,\K f8 folgt 7.\,\D f7\#, und
	6\,\ldots\,\D b3+ 7.\,\D xb3+ verzögert das Matt nur.}

\Schritt{7w}{straightmove}{e4}{e8}{e1-e8}
{7.\,\T e8+! \quad \T xe8}
{Das zweite Turmopfer. Die Dame ist abgeschnitten, also muss der Turm von e1 schlagen. Genau das
	wollte Weiß: Jetzt steht ein schwarzer Turm auf der 8.~Reihe, der später auf g8 schlagen muss.
	Außerdem versperrt er nun selbst der Dame den Weg nach g8.}

\Schritt{8w}{knightmove}{d8}{f7}{h8-g8}
{8.\,\Sp f7+ \quad \K g8}
{Zum dritten Mal Springerschach auf f7, zum dritten Mal bleibt nur \K g8.}

\Schritt{9w}{knightmove}{f7}{h6}{g8-h8}
{9.\,\Sp h6+ \quad \K h8}
{\textbf{Doppelschach:} Springer und Dame geben gleichzeitig Schach. Gegen ein Doppelschach hilft nur
	ein Königszug, schlagen (g7xh6) oder dazwischenziehen geht nicht. Nach 9\,\ldots\,\K f8 wäre
	10.\,\D f7 matt, also wieder \K h8.}

\Schritt{10w}{straightmove}{a2}{g8}{e8-g8}
{10.\,\D g8+!! \quad \T xg8}
{Das Damenopfer. Der König darf die Dame nicht schlagen, weil der Springer von h6 aus g8 deckt.
	Der Turm muss nehmen und mauert damit seinen eigenen König ein.}

\Schritt{11w}{knightmove}{h6}{f7}{}
{11.\,\Sp f7\#}
{\textbf{Ersticktes Matt.} Der König steht im Schach und hat kein Feld: g8, g7 und h7 sind von
	eigenen Figuren besetzt. Der Springer kann nicht geschlagen werden, denn der Turm auf g8 deckt f7
	nicht und die Dame auf b8 kommt nicht heran. Weiß hat zwei Türme und die Dame geopfert und setzt
	mit dem letzten Springer matt.}

\end{document}
